%% Chapter 10, Section 10.4 — Exercises
%% Source: A First Course in Monte Carlo Methods, Sanz-Alonso & Al-Ghattas
%% Book pages 123–126

\section{Exercises}
\label{sec:10.4}

\begin{enumerate}

\item[10.1] Consider a univariate Gaussian with unknown mean $\mu$ and variance $\tau^{-1}$. For
convenience, we work in terms of the precision $\tau$, the reciprocal of the variance. We are
provided with samples $y^{(1)}, \ldots, y^{(K)}$ and the setting of the problem is as follows:
\begin{align*}
  \mu &\sim \mathcal{N}(\mu_0,\, (\lambda_0 \tau)^{-1}), \\
  \tau &\sim \operatorname{Gamma}(a_0,\, b_0), \\
  y^{(1)}, \ldots, y^{(K)} &\overset{\mathrm{i.i.d.}}{\sim} \mathcal{N}(\mu,\, \tau^{-1}).
\end{align*}
Here, $\mu_0$, $\lambda_0$, $a_0$, and $b_0$ are hyperparameters that describe the prior
distribution on our unknown parameters $\theta = (\mu, \tau)$. The mean-field approximation gives
\[
  g(\mu, \tau) \coloneqq g_\mu(\mu)\, g_\tau(\tau).
\]

Show that
\begin{align*}
  g_\mu^*(\mu) &\sim \mathcal{N}(\mu_K,\, \tau_K^{-1}), \\[4pt]
  \mu_K &= \frac{\lambda_0 \mu_0 + K\bar{y}}{\lambda_0 + K}, \\[4pt]
  \tau_K &= (\lambda_0 + K)\,\mathbb{E}_\tau[\tau],
\end{align*}
where $\bar{y} = \frac{1}{K}\sum_{k=1}^K y^{(k)}$ denotes the sample mean. Likewise, show that
\begin{align*}
  g_\tau^*(\tau) &\sim \operatorname{Gamma}(a_K,\, b_K), \\[4pt]
  a_K &= a_0 + \frac{K+1}{2}, \\[4pt]
  b_K &= b_0 + \frac{1}{2}\,\mathbb{E}_\mu\!\left[\sum_{i=1}^{K}(y_i - \mu)^2
         + \lambda_0(\mu - \mu_0)^2\right].
\end{align*}

The distribution of $\mu$ is determined by $\mu_K$ and $\tau_K$, whereas the distribution of $\tau$
is determined by $a_K$ and $b_K$. Notice that $a_K$ and $\mu_K$ can both be calculated from the
data and known hyperparameters. However, $b_K$ depends on $\mu_K$ and $\tau_K$, and $\tau_K$
depends on $a_K$ and $b_K$. So there is a circular dependence between $\tau_K$ and $b_K$, which
leads to the following CAVI algorithm:

\begin{algorithm}[H]
  \caption{CAVI for Univariate Gaussian}
  \label{alg:10.3}
  \begin{algorithmic}[1]
    \Require Data $y$, hyperparameters $a_0$, $b_0$, $\mu_0$, $\lambda_0$.
    \State Calculate $\mu_K$ and $a_K$ using data and hyperparameters.
    \State Initialize $b_K$.
    \State Use current $b_K$ to calculate $\tau_K$.
    \State Use current $\tau_K$ to calculate $b_K$.
    \State Repeat steps 3 and 4 until convergence.
  \end{algorithmic}
\end{algorithm}

Illustrate the performance of the method. To that end:
\begin{enumerate}
  \item Generate data $y$ from the model with true parameter values of your choice.
  \item Choose hyperparameter values.
  \item Plot level curves of the variational distributions you obtain along with the contour plots
        of the true posterior.
\end{enumerate}

\item[10.2] Consider again the original EM example by Dempster et al.\ 1977. Suppose that data
$(125, 18, 20, 34)$ are collected.

\begin{enumerate}
  \item[\textit{a)}] Construct an EM algorithm to find the maximum likelihood estimator of $\theta$.
  \item[\textit{b)}] Construct a Monte Carlo EM algorithm to find the maximum likelihood estimator
        of $\theta$. Compare your results with those of the algorithm designed in part a) using a
        Monte Carlo experiment evaluating the variability of the Monte Carlo EM sequence.
\end{enumerate}

\item[10.3] We say that a random variable $Y$ is distributed according to a Gaussian mixture model
with $C$ clusters if
\begin{equation}
  Y \sim \sum_{c=1}^{C} \pi_c\, \mathcal{N}(\mu_c, 1).
  \label{eq:10.6}
\end{equation}
For $1 \leq c \leq C$, $\mu_c$ represents the mean of the $c$-th cluster and $\pi_c$ the
probability with which a draw from the mixture model belongs to the $c$-th cluster. Thus, it holds
that $0 < \pi_c < 1$ and $\sum_c \pi_c = 1$. In this exercise, you will use an EM algorithm to
estimate the vector $\mu = (\mu_1, \ldots, \mu_C)$ of cluster means and the vector
$\pi = (\pi_1, \ldots, \pi_C)$ of mixing coefficients given data $y = \{y^{(1)}, \ldots, y^{(K)}\}$
independently drawn from the Gaussian mixture model. The number $C$ of clusters will be assumed to
be known.

\begin{enumerate}
  \item[\textbf{1.}] \textbf{Setting and Preparation}

  Consider a Gaussian mixture model with $C = 3$ clusters. Sample $K = 10^4$ draws
  $y^{(1)}, \ldots, y^{(K)}$ from the Gaussian mixture model with true parameters
  $\mu = (0, -1, 1)$ and $\pi = (1/4, 1/4, 1/2)$. Your goal will be to recover these parameters
  from the simulated data. To obtain each draw $y^{(k)}$, $1 \leq k \leq K$, do the following:

  \begin{enumerate}
    \item[\textit{a)}] Sample a latent variable $Z^{(k)}$ with p.m.f.\ $\mathbb{P}(Z^{(k)} = c) = \pi_c$,
          $1 \leq c \leq 3$. Note that to generate the simulated data you should use the true mixing
          coefficients $\pi = (1/4, 1/4, 1/2)$. The random outcome $c \in \{1, 2, 3\}$ of the draw
          $Z^{(k)}$ will be interpreted as the cluster assignment of the $k$-th data point $y^{(k)}$.
          We denote $Z = (Z^{(1)}, \ldots, Z^{(K)})$.
    \item[\textit{b)}] Given that $Z^{(k)} = c$, draw $y^{(k)} \sim \mathcal{N}(\mu_c, 1)$. Note
          that to generate the simulated data you should use the true cluster means vector
          $\mu = (0, -1, 1)$.
  \end{enumerate}

  Plot a histogram of all draws $y^{(k)}$ you simulated in (b) using three different colors
  according to the value of the corresponding latent variable $Z^{(k)}$.

  \item[\textbf{2.}] \textbf{EM Algorithm}

  \begin{enumerate}
    \item[\textit{c)}] Prove that the full log-likelihood admits the following factorization:
    \[
      \log f(y, z|\mu, \pi) = \log f(y|z, \mu) + \log f(z|\pi).
    \]

    \item[\textit{d)}] Prove that
    \[
      \mathbb{E}_{Z \sim f(z|y,\mu^{(\ell)},\pi^{(\ell)})}\!\bigl[\log f(y, Z|\mu, \pi)\bigr]
      = \sum_{k=1}^{K}\sum_{c=1}^{C} w_{kc}^{(\ell)}
        \bigl(\log \mathcal{N}(y^{(k)};\, \mu_c, 1) + \log \pi_c\bigr),
    \]
    where
    \[
      w_{kc}^{(\ell)} \coloneqq \mathbb{P}\!\bigl(Z^{(k)} = c\,|\,y^{(k)},\, \mu^{(\ell)},\, \pi^{(\ell)}\bigr)
      = \frac{\pi_c^{(\ell)}\,\mathcal{N}(y^{(k)};\,\mu_c^{(\ell)},\,1)}
             {\sum_{c'=1}^{C}\pi_{c'}^{(\ell)}\,\mathcal{N}(y^{(k)};\,\mu_{c'}^{(\ell)},\,1)},
    \]
    and $\mathcal{N}(y^{(k)};\,\mu_c^{(\ell)},\,1)$ denotes a $\mathcal{N}(\mu_c^{(\ell)},\,1)$
    p.d.f.\ evaluated at $y^{(k)}$.

    \item[\textit{e)}] Prove that $\mu^{(\ell+1)}$ and $\pi^{(\ell+1)}$ given below maximize the
    expectation:
    \[
      \pi_c^{(\ell+1)} = \frac{\sum_{k=1}^{K} w_{kc}^{(\ell)}}{K}, \qquad 1 \leq c \leq C,
    \]
    \[
      \mu_c^{(\ell+1)} = \frac{\sum_{k=1}^{K} w_{kc}^{(\ell)}\, y^{(k)}}{\sum_{k=1}^{K} w_{kc}^{(\ell)}},
      \qquad 1 \leq c \leq C.
    \]
    \textit{Hint}: note that in the maximization you should enforce the constraint
    $\sum_c \pi_c = 1$.

    \item[\textit{f)}] Implement an EM algorithm using your simulated data and the calculations
    above. Plot the estimated parameters against the iterations, comparing with their true values.

    \item[\textit{g)}] Derive the formula for the log-likelihood, and plot the value of the
    log-likelihood against the iterations.
  \end{enumerate}
\end{enumerate}

\item[10.4] Consider a Gaussian mixture model with unknown cluster variances
\[
  Y \sim \sum_{c=1}^{C} \pi_c\, \mathcal{N}(\mu_c,\, \sigma_c^2).
\]
Following a similar strategy as in the previous exercise, design an EM algorithm to estimate the
unknown parameters $\mu = (\mu_1, \ldots, \mu_C)$, $\pi = (\pi_1, \ldots, \pi_C)$, and
$\sigma = (\sigma_1, \ldots, \sigma_C)$. Implement your approach using simulated data with true
parameters $\mu = (-2, 0, 2)$, $\sigma^2 = (0.01, 1, 0.16)$, and
$\pi = \bigl(\tfrac{1}{6}, \tfrac{1}{6}, \tfrac{2}{3}\bigr)$. The number $C = 3$ of clusters is
assumed to be known.

\begin{enumerate}
  \item[\textit{a)}] Derive the formulas for the E and M steps.
  \item[\textit{b)}] Plot the estimated parameters against the iterations, comparing with their
  true values.
  \item[\textit{c)}] Derive the formula for the log-likelihood, and plot the value of the
  log-likelihood against the iterations.
\end{enumerate}

\end{enumerate}
